% Proposed insertion only. No master/include/source selection is changed.
% Inherit the selected master's Noto Sans/Fira Math setup; no new font defaults.
\section{More data must survive a harder test}

A thousand repetitions of one reach still describe one kind of reach. A movement model trained across people, recording days, tasks and bodies could encounter relationships that a narrower collection never contains. Whether that breadth improves a particular prediction is the scaling hypothesis proposed here. Increasing the amount of qualified data and increasing its diversity are separate interventions. Neither guarantees improvement, and the proposed Large Movement Model has not been trained or validated as a body-wide world model.

\subsection{Four source groups, distinct observations}

The proposed training programme draws on at least four groups: neural observations; muscle and other myographic observations; measured motion and displacement; and data generated by physics-based motion or musculoskeletal simulation. These are source groups, rather than a claim of exhaustive coverage or four fully acquired corpora. The neural scope includes suitable brain, spinal, peripheral and relevant visceral observations, with their recording sites, modalities and task roles stated separately. The present emphasis is motor--sensory. Myoception is the programme's term for the sensory side of musculoskeletal contraction and movement, including proprioception, tactile feedback and other task-relevant sensations linked to action. It marks proposed programme scope, without establishing a new external scientific standard or a body-wide simultaneous recording.

Tactile stimulation, proprioceptive feedback, contact, force and task context require different observation relationships. A known tactile stimulus records what was delivered; a contact sensor records its specified contact event; a force instrument measures its calibrated physical quantity. Kinematics may help estimate body configuration without directly recording sensory afferents or the subject's proprioceptive experience. EMG supplies muscle electrical activity, not direct peripheral nerve or sensory recording. Measured stimuli, reported sensation, recorded neural features and model estimates must keep their identities even when they describe one action.

The fourth group contributes training examples as well as physical constraints. A qualified generator could supply simulated trajectories, state transitions, muscle--tendon or contact-force outputs and declared movement priors under specified controls. Every example needs its simulator and body-model version, parameters, solver, control representation and source lineage. A joint-target controller and a muscle-driven controller expose different action spaces; one must not be reported as the other. Motion retargeted from a recording retains that recording's lineage and the transformation, while simulated states remain simulated. A generated force is not an independent force measurement.

Before using such examples, calibrate the generator's relevant parameters against suitable measured support, retaining calibration uncertainty and training-only role boundaries. Compare a fixed-real-data reference with a simulation-augmented arm under matched permitted information, compute and original criteria. Record real-example exposure, simulated exposure and updates in each arm: a fixed real corpus alone does not match training exposure. Report simulated amount, simulated diversity and real acquisition amount separately; then test the synthetic contribution against held-out real observations. Keep those real tests out of generator tuning, calibration, reward choice and model selection. Errors in contact, dynamics or source correspondence must remain ways for the synthetic contribution to fail. Generated data cannot supply independent biological ground truth for its own assumptions. The generator's methods, implementation and use conditions still need checking; no frontier implementation or dataset is selected here.

The four source groups and five component roles describe different axes. Brain signals, muscle signals, measured movement, a virtual body and a learning model describe jobs in the proposed system. A virtual body can both generate training examples and constrain forecasts; the learning model can use several source groups. Neither arrangement closes the inventory to other useful, properly measured information.

\subsection{Separate amount from diversity}

For each modality, training amount must identify what was counted: original acquisition groups, people or animals, recording duration, trials, observed target values and distinct training windows. Overlapping windows and repeated passes do not create independent acquisitions. The study must also report diversity as a description of the training distribution: organism, subject, session, task, body region, recording modality, apparatus, environment and measurement quality. One additional subject and one additional sensor are different changes. A count of datasets cannot make them interchangeable.

The proposed amount comparison would retain a declared source distribution while increasing eligible training support. Its diversity comparison would retain a declared amount while changing specified dimensions of that distribution. A crossed design could estimate their interaction. All comparisons would keep the architecture, allowed input information, target, preprocessing and model-selection rules fixed. Training exposure, repetitions, updates and measured computational cost must be reported together: equal update counts do not by themselves establish equal compute or equal information. Where different modalities or recording quality prevent a matched comparison, that limitation becomes part of the result rather than a hidden explanation for it.

Each new comparison needs its own prospectively frozen roster, budgets, hypotheses, seeds, controls and decision rules. It does not replace the adopted continuous WAY-EEG-GAL development design or alter its fixed source/target roles, updates or selection. Nor may it reopen the reserved LINK dates, ST or other protected tests. Training-only calibration and normalization follow original acquisition groups through every trim, derivative, alias and window. An exposed recording cannot become a fresh test by changing its name.

Generalization must name the change being tested. An excluded session tests a different question from an excluded subject, task, apparatus, dataset or species. The proposed evaluation would report these domains separately, with a fixed target-only/no-transfer reference, the appropriate incorrect-correspondence control, simpler predictors and the original null, shuffle, confound and causal-time controls. New data may help familiar cases while harming an unfamiliar body or task; those failures and any negative transfer must remain visible. A pooled improvement cannot overrule an original no-harm requirement.

\subsection{A power law is a candidate description}

For one defined error measure, one evaluation population and one fixed diversity condition, a candidate relationship is
\[
E(N)=E_{\infty}+A\left(N/N_0\right)^{-\alpha}.
\]
Here $N>0$ is the registered measure of training amount, $N_0>0$ is a fixed reference amount in the same units, and $E(N)$ is a specified nonnegative held-out error. For this decreasing-error candidate, require $E_{\infty}\geq0$, $A\geq0$ and $\alpha>0$. $E_{\infty}$ and $A$ have the error's units; the ratio and exponent are dimensionless. With $A>0$ the curve decreases strictly as $N$ increases; $A=0$ gives constant error. A positive exponent alone is insufficient. Any upper bound of the chosen metric must also be respected throughout the declared fitting and prediction range. The fitted floor would be a parameter of this candidate, not an independently established irreducible biological error. This form would establish neither a universal law, unlimited extrapolation nor improvement from every additional dataset.

The form is a hypothesis, with no fitted exponent reported here. Its decreasing family does not constrain what the observations are allowed to show. Constant performance, saturation, increasing error, non-monotonic change and domain-specific deterioration must be credible alternatives, and negative transfer must be reported even when this curve cannot describe it. The fitting range, alternative models, uncertainty procedure and predictive check at withheld scaling levels must be fixed before inspecting the relevant outcomes. Related windows cannot serve as independent uncertainty replicates. An apparent curve driven by one additional subject, a changed target definition or a less demanding test distribution must be diagnosed at that level.

The errors themselves retain their meanings. A normalized EMG-target error, an image-pixel displacement error and a physical joint-angle error cannot be pooled into a biological accuracy percentage without an independently justified scoring rule. Report the original units, per-domain errors, sample denominators, missingness and uncertainty alongside any declared dimensionless summary. Acquisition diversity, computational model size and statistical averaging remain separate axes.

\subsection{Complementary recordings leave missing links}

FALCON M1-A offers a concrete neural--muscle candidate: the fixed DANDI release describes simultaneous M1 activity and preprocessed intramuscular EMG from one rhesus macaque. Its three split labels do not represent three animals. No released synchronized motion-capture trajectory is established by the reviewed metadata. The published target processing is acausal, so an offline result using that target would not establish a purely causal real-time system. Its small minival subset comes from the calibration data and is used to check submission format; it is not an independent validation population. These constraints belong beside any proposed use of the dataset. \href{https://doi.org/10.48324/dandi.000941/0.241029.1405}{FALCON M1-A release}; \href{https://papers.nips.cc/paper_files/paper/2024/file/8c2e6bb15be1894b8fb4e0f9bcad1739-Paper-Datasets_and_Benchmarks_Track.pdf}{FALCON methods}.

AJILE12 offers a different relationship: intracranial recordings and video-estimated, two-dimensional upper-body pose across twelve clinical participants. Its pose coordinates are image pixels, not instrumented three-dimensional motion capture; paired EMG is not established in the reviewed release. Pose-label construction is participant-adapted, and possible exposure of a particular future held-out day to the pose estimator remains unverified. The numerical alignment uncertainty also remains to be qualified. A neural predictor's unseen subject or day therefore does not automatically make its target-label construction independent. \href{https://doi.org/10.1038/s41597-022-01280-y}{Peterson et al., 2022}; \href{https://doi.org/10.48324/dandi.000055/0.220127.0436}{AJILE12 release}.

These sources provide complementary candidate links. A macaque's neural--EMG example and a person's neural--pose example never become one simultaneous neural--muscle--movement observation. Qualified unpaired pretraining could learn source-specific representations; an explicit transfer model could then be tested against actual target measurements. It would still have to earn its cross-domain correspondence. These release descriptions establish possible source relationships, rather than downloaded recordings or a trained mapping. The exact use conditions, file integrity, clocks, units, earlier exposure and independence of the chosen partitions still need checking before training.

\section{A world model must account for the body it moves}

Consider a hand closing on an object. The exterior changes: fingers displace, contact begins and the object may move. Inside the model, muscle--tendon forces and joint motion must account for that change under the declared load. The proposed world model joins those questions. Its interior initially means the specified musculoskeletal states and mechanics; its exterior means the supported body motion, objects, contacts and environment. The phrase does not imply that every organ, neural pathway or private sensation has been reconstructed.

NeuroTranslate would supply conditional mappings between declared observations and targets. MoCap ReComp names the proposed source-preserving composite movement reference; a video-derived pose collection does not thereby become that instrument. MyoSim Replay names the proposed physical simulation component. Their proposed combination is the Large Movement Model. A neural composite here would retain identified recording types and source-conditioned representations. Combining them in software would not identify equivalent neurons across people or species, or assemble a measured whole nervous system.

\subsection{Keep observations and internal state distinct}

A prospective model could represent mechanical state by joint configuration $q$, generalized velocity $v$, muscle-model activation $a$ and declared environmental state $c$. For a rotational joint, $q$ and $v$ may use radians and radians per second; translations require their own length units. The simulator's state dimensions and conventions determine the actual representation. A learned transition would predict a distribution over later states, conditioned on the permitted past observations, task context and any declared control input. At inference, withheld measurements and future samples remain unavailable even if they supplied training targets.

Each sensor needs a separate observation relationship. A neural recording measures its named electrical or imaging quantity. EMG measures muscle electrical activity; its processed amplitude is not automatically calibrated muscle activation or force. Tracking supplies observed or estimated geometry with its calibration and uncertainty. A simulated muscle force remains a model output unless an independent force reference tests it. Missing observations receive an availability mask, not a fabricated zero or a synthetic measurement label. The reverse prediction of recorded neural features remains conditional and potentially non-unique.

For a declared articulated-body model, one proposed mechanical check can be written schematically as
\[
M(q)\dot v+h(q,v)
 =R(q)F_{\mathrm{mt}}+J_c(q)^{\mathsf T}f_c+\tau_{\mathrm{ext}}.
\]
$M$ is the model's inertia operator and $h$ collects its declared velocity and gravitational terms. Muscle--tendon forces $F_{\mathrm{mt}}$ act through the moment-arm or generalized-force mapping $R$; contact forces $f_c$ act through contact Jacobian $J_c$; $\tau_{\mathrm{ext}}$ contains other declared generalized forces. Each equation component must use the model's compatible coordinates and generalized-force units, with the coordinate-to-velocity mapping stated where $v\neq\dot q$. This is a proposed consistency relation for a specified simulator, not an inferred equation of the participant's complete biology. Geometry, inertia, muscle parameters, external loads and contact assumptions still require independent support. OpenSim provides established musculoskeletal modelling methods, rather than validation of this proposed assembly. \href{https://journals.plos.org/ploscompbiol/article?id=10.1371/journal.pcbi.1006223}{Seth et al., 2018}.

Predicted movement must be compared in aligned time and coordinates with independent movement references. Internal force claims need appropriate force measurements; matching the visible path alone cannot identify the force distribution that produced it. Test contacts, joint limits, energy accounting, gravity, friction and solver sensitivity within the claimed domain. A model that exploits a numerical correction or receives privileged future targets can improve its score without improving the intended forecast. Report that failure as such. A more lifelike rendering cannot repair it.

The first useful result may be a limited forecast with explicit missing state, rather than a complete body. Its tests would isolate the neural representation, muscle observations, movement prior and physics constraint through matched ablations and external outcomes. Any later policy refinement would need suitable observations of actions, their consequences and rewards, with its reward calculation checked against those independent outcomes. Predictive fitting, offline replay, simulated control, physical intervention and retained biological learning remain different experiments.

% R4 insertion candidate only; no existing manuscript is modified.
% Inherit Noto Sans/Noto Sans Mono/Fira Math from the selected master.
% Requires the shared result table and the proposed supplementary bibliography.
\subsection{One muscle pulse, two simulated trajectories}

A muscle command should change the body the simulator computes. We tested that limited relationship in one installed human right-hand mechanism, \texttt{myohand\_r}, attached to a fixed passive torso. The question was whether a specified excitation changes the recorded mechanical state under a matched comparison. No neural recording, measured EMG or measured movement supplied the input or reference. This is a numerical software diagnostic within the wider mechanics proposal; it does not validate a biological hand or select a training generator. The model came from \texttt{myo\_sim} 0.2.3 at commit \texttt{93b0ca8f4ec90c9899ee7f05fee561e9911da91b}, with MuJoCo 3.14.0 and NumPy 2.5.3.\cite{myohandMyoSimPinned2026,myohandMuJoCo3140,mujoco2012,myohandNumericalRecord20261105}

\subsubsection{Mechanism and matched comparison}

The compiled mechanism has 23 scalar hinge degrees of freedom, 39 muscle--tendon actuators, 39 activation states and zero sensors. Its coordinates are joint angles in radians. Both branches began from matching full recorded initial states, including joint positions, velocities and activation states. The control branch received no additional excitation pulse; the other received dimensionless excitation 0.2 at the actuator named \texttt{FDS2\_r}, with all other controls zero. This input is a muscle-model command, distinct from a joint-position target or an electrical recording.

Both branches used the native Euler setting with $\Delta t=0.002\,\mathrm{s}$ for 100 explicit steps, giving a nominal horizon $T=0.200\,\mathrm{s}$ per branch. Pulse steps 10 through 59, counted from zero, use the pre-step interval $0.020\leq t<0.120\,\mathrm{s}$. Gravity and passive dynamics remain active in the comparison: zero added excitation need not leave the hand still. Initial and final states were each synchronized with an explicit forward evaluation. The record therefore contains 200 explicit integration steps and four explicit forward evaluations. Two compilation calls and one internal construction-time forward evaluation are inferred from the executed source path; these are not counts of every internal native solver operation.

Construction also matters. Although the optional \texttt{inertia\_floor} argument remained \texttt{None}, the native builder changed each diagonal inertia component of four fixed scaffold bodies from $1$ to $10^{-4}\,\mathrm{kg\,m^2}$. That transformation belongs to the tested mechanism and must remain with its provenance. No additional physical-parameter retuning or fallback to \texttt{implicitfast} was used. Position, orientation, model scale, contact behaviour and biological fidelity have not been independently calibrated by this diagnostic.

\subsubsection{Recorded response}

Let $P$ denote the pulse branch and $C$ the matched control. At the synchronized final state, define
\[
\Delta q_{\infty}=\max_{1\leq j\leq23}|q_{P,j}(T)-q_{C,j}(T)|,
\qquad
\Delta a=a_{P,\mathrm{FDS2}}(T)-a_{C,\mathrm{FDS2}}(T),
\]
\[
\Delta F=F_{P,\mathrm{FDS2}}(T)-F_{C,\mathrm{FDS2}}(T).
\]
Here $a$ is the selected model activation, with no physical unit, and $F$ is the selected signed tendon-actuator force in newtons under the model's convention. The coordinate maximum is across the 23 hinges at the final time, rather than across the whole trajectory. The force uses the final synchronized forward evaluation; a force retained from the last integration step need not refer to that same state.

% Shared by all four proposed R3 inserts; no added packages or font defaults.
\begin{table}[htbp]
\centering
\caption{Final contrasts for one simulated right-hand mechanism. All values compare the pulse branch with its matched control at nominal $T=0.200\,\mathrm{s}$. The angle value is the largest absolute separation across 23 hinges at that endpoint; activation and force retain the pulse-minus-control sign. Activation and force concern the selected \texttt{FDS2\_r} actuator. Force uses the synchronized final forward evaluations and the model's sign convention. These are simulated contrasts, not errors against biological measurements. The displayed digits preserve the recorded numerical values, without establishing biological accuracy.\cite{myohandNumericalRecord20261105}}
\label{tab:myohand-numerical-response}
\begin{tabular}{p{0.53\linewidth}p{0.26\linewidth}p{0.10\linewidth}}
\hline
Quantity & Contrast & Unit \\
\hline
Largest absolute final joint-coordinate separation, $\Delta q_{\infty}$ & $0.5027347840656727$ & rad \\
Selected final activation difference, $\Delta a$ & $0.05205729904534248$ & none \\
Selected final tendon-actuator force difference, $\Delta F$ & $-6.407722136774943$ & N \\
\hline
\end{tabular}
\end{table}


An independent standard-library audit of the recorded evidence completed 132 checks without a failed check and without rerunning the model. The required recorded numerical channels were finite, and recorded native warning counts were zero. The four stored endpoint mass-matrix spectra contained positive values. The raw mass matrices and residual operands were not independently exported, so the audit checked those stored diagnostics rather than solving their underlying matrices again. No contact was observed in this case. At the final control endpoint, the recorded soft-limit excess was $0.0004121829780273552\,\mathrm{rad}$; these soft constraints must not be described as exact hard joint bounds. The diagnostic gives no conservation proof, convergence study or physiological magnitude threshold.

The nonzero contrasts establish an input-to-state response in this configuration. They are differences between two simulated trajectories, not errors against human measurements. The 200 dependent integration steps are neither 200 subjects nor 200 independent experiments; no statistical $p$-value is inferred. Extra decimal places preserve the numerical record without establishing biological accuracy or repeatability across cases.

\subsubsection{The next claim needs a different reference}

The run used a CPU dynamics route. Readback showed a one-thread BLAS pool and a zero-sized dynamics pool, while transient compiler parallelism remained unmeasured. The 263 observed loaded-file identities form a bounded record, not a complete system closure. Resource and performance qualification remain unresolved: process high-water RSS was reported as 5,187,464 KiB and can retain pre-exec history, so it cannot be assigned to the hand model or used to establish a model-only 3 GiB bound. Tool/native timing near 448 ms and in-probe timing near 491.927 ms have unreconciled boundaries. Neither model RAM use, complete CPU/GPU partitioning, useful end-to-end latency nor a speedup follows from these values.

This mechanism supplies simulated activation, force and motion; it has no sensory measurements. Myoception---the programme's proposed sensory side of muscle-driven action---would require its own task-relevant sensory evidence. No afferent recording, learned controller, biological feedback loop, trained Large Movement Model or transfer result is demonstrated. The four proposed training-source groups remain neural observations, myography, measured motion and physics-generated examples. Before this generator could contribute training data, its relevant mechanics, controls, calibration, uncertainty, source rights and training lineage would need qualification, followed by independent held-out real tests. The amount/diversity hypothesis remains open. No protected test or quarantined source is reopened by this diagnostic.

The current record supplies no successful scaling curve or trained body-wide world model. LINK's eighteen fits and 720 saved epochs still fail both model-admission gates and all five benefit claims. New data would have to improve a new, properly qualified comparison; it cannot revise that negative result after the fact. The proposed programme becomes stronger when it names the particular relationship that survives a harder test, and equally when it identifies the one that fails.
